% This nuaa-AMS.tex is a latex-beamer template using the beamer theme for the university of Amsterdam produced by Nico Schlömer, modified by Felix Ding.
% URL: https://github.com/nschloe/ua-beamer
% 
% Copyright 2009--2011 by Nico Schlömer
% This file may be distributed and/or modified
% 1. under the LaTeX Project Public License and/or
% 2. under the GNU Public License.

% Required files to compile by pdflatex:
% 1. beamerthemeUniversiteitAntwerpen.sty
% 2. beamercolorthemeUniversiteitAntwerpen.sty
% 3. beamerfontthemeUniversiteitAntwerpen.sty
% 4. beamerinnerthemeUniversiteitAntwerpen.sty
% 5. beamerouterthemeUniversiteitAntwerpen.sty
% 6. bibliography_file.bib
% 7. nuaa_logo.pdf
% 8. nuaa_name.pdf
% 9. uabottomwave.pdf is required only for light theme
% 10. uaBackgroundDarkBlue.pdf is required only for dark theme
% 11. uaBackgroundLightBlue.pdf is required only for dark theme

% \documentclass[handout]{beamer}
\documentclass[CJKutf8,compress]{beamer}

% switch the following lines to choose light and dark themes
\usetheme[light,framenumber,totalframenumber]{UniversiteitAntwerpen}
%\usetheme[dark]{UniversiteitAntwerpen}
\usepackage{CJKutf8}
\usepackage{libertine}

\begin{document}
\begin{CJK*}{UTF8}{gkai} % support for Chinese

  \title{Main Title 主标题}
  \subtitle{Subtitle 副标题} % optional
  \author{Author 作者}

  \begin{frame}
    \titlepage
  \end{frame} % ================================================================
  
  \begin{frame}{Table of Contents 目录}
    \tableofcontents
  \end{frame} % ================================================================
  
  \section{1. Introduction}
  \label{Sec:introduction}
  \begin{frame}{1. Introduction 引言}
    \framesubtitle{frame subtitle 页面副标题}
    This nuaa-AMS.tex is a \LaTeX \ Beamer template using the Beamer theme for the university of Amsterdam\cite{AMSTheme} produced by Nico Schlömer.

    \bigskip

    此nuaa-AMS.tex文件是一个使用了Nico Schlömer制作的阿姆斯特丹大学主题\cite{AMSTheme}的\LaTeX \ Beamer模板。

    \begin{figure}
      \begin{center}
        \includegraphics[scale=0.1]{latex.png}
      \end{center}
      \caption{\LaTeX \ beamer}
      \label{Fig:latex_beamer}
    \end{figure}
  \end{frame} % ================================================================

  \section{2. Mathematical model}
  \label{Sec:model}
  \begin{frame}{2. Mathematical model 数学模型}
    Lists 列表

    \begin{itemize}
    \item Apple
    \item Orange
    \item Banana
    \end{itemize}

    \begin{enumerate}
    \item Monday
    \item Tuesday
    \item Wednesday
    \end{enumerate}

    \begin{description}
    \item[Description list] is a type of list to describe items.
    \item[Description list] 是一种用于描述的列表。
    \end{description}
  \end{frame} % ================================================================
  
  \begin{frame}{2. Mathematical model 数学模型}
    \alert{$\backslash$pause命令}可用于逐次展示itemize、enumerate中的条目。

    \begin{itemize}
    \pause
    \item Apple
    \pause
    \item Orange
    \pause
    \item Banana
    \end{itemize}
  \end{frame} % ================================================================
  
  \section{3. Empirical experiments}
  \label{Sec:experiments}
  \begin{frame}{3. Empirical experiments 实验}
    Blocks 区块

      \begin{definition}[Pythagorean theorem]
        The Pythagorean theorem is a fundamental relation in Euclidean geometry among the three sides of a right triangle.
      \end{definition}

      \begin{theorem}[Pythagorean theorem]
        $a^2 + b^2 = c^2$
      \end{theorem}
  \end{frame} % ================================================================

  \begin{frame}{3. Empirical experiments 实验}
    Blocks 区块

      \begin{exampleblock}{For example}
        $3^2 + 4^2 = 5^2$
      \end{exampleblock}

      \begin{alertblock}{Note}
         Note that the Pythagorean theorem can only be applied to right triangles.
      \end{alertblock}
  \end{frame} % ================================================================

  \begin{frame}
    \begin{center}
      \Huge{\bf{Thank you!}}
    \end{center}
  \end{frame} % ================================================================

  \bibliographystyle{IEEEtran} % comment this line if nothing is cited
  \bibliography{bibliography_file} % comment this line if nothing is cited

\end{CJK*}
\end{document}